% s-163-p163-m001; original PDF 163
\Aphorism{s-163-p163-m001}{248}
In geistigen Dingen ist Altruismus, Apostelthum, Parteiwesen höchst verdächtig. Wer nur einmal einen Blick in die Fülle des zu Geniessenden gethan hat, wird immer mehr die Lust verlieren, auf einer Zwischenstufe Halt zu machen, um erst die „Anderen“ nachzuholen und Anhänger zu werben. Mag Jeder doch sich selbst forthelfen; wer heute Hunger hat und nichts zu essen findet, ist ein unverbesserlicher Dummkopf. Unser Zeitalter ist ja intellectuell gewiss nicht, was es vielleicht wirthschaftlich sein mag: ein Zeitalter der Nothstände und der ungerechten Gütervertheilung — es giebt keinen Leckerbissen indischer, griechischer oder sonstwelcher Herkunft, den nicht Jeder mit geringem Aufwand erstehen könnte. Es ist wirklich nicht recht zweck- und zeitgemäss, von der Bildung und Erziehung Anderer so viel Wesens zu machen oder gar die eigene darüber zu versäumen.
\par

